\chapter{Chygraphs}
\label{ch:chygraphs}

\section{A paper is not a set of links}
\label{sec:paper}

In this chapter we define the chygraph and show what it represents that a graph
or a hypergraph cannot. We motivate it with the scientific literature, which
has been mapped many times over, whose maps disagree with one another, and
where the reason they disagree is exactly the subject of the chapter.

The system itself is simple. A paper is written by some authors. It cites
some other papers. Those papers were written by their own authors and cited
their own predecessors. That is the whole system --- documents, the people who
wrote them, and the pointers between documents --- and it has been drawn as a
network in at least four different ways \citep{newman2001,barabasi2002}.

The \emph{citation network} keeps the papers and the pointers: a node is a
paper, a directed link is a citation \citep{redner1998}. It is the right object
if the question is how knowledge flows, and it has nothing to say about who
wrote anything. The \emph{co-authorship network} keeps the people: a node is an
author, and two authors are linked if they have written a paper together. It is
the right object if the question is who works with whom, and it has nothing to
say about citation. The \emph{bipartite author--paper graph} keeps both kinds of
thing, authors and papers, with a link wherever a person wrote a document; it
loses the citations again. And the \emph{author hypergraph} replaces the
co-authorship links by hyperedges, one per paper, so that a five-author paper is
one object rather than ten pairwise links \citep{vasilyeva2021}.

Four maps, four different objects, and each of them made by throwing away
something that was in the system to begin with (Fig.~\ref{fig:projections}).
They also throw away different things, so a result proved on one of them
cannot be quoted on another. That is what happens when a system whose parts
contain other parts is forced into a language
whose only relation is ``these two things are joined.''

\begin{figure}[t]
\centering
\begin{adjustbox}{max width=\linewidth}
\begin{tikzpicture}
% ---- (a) citation network
\begin{scope}
  \node[tb,anchor=west] at (-0.1,2.15) {(a) citation network};
  \node[cpx] (p1) at (0.35,1.35) {}; \node[cpx] (p2) at (1.55,1.50) {};
  \node[cpx] (p3) at (1.00,0.65) {}; \node[cpx] (p4) at (2.20,0.60) {};
  \draw[incd] (p1)--(p2) (p3)--(p1) (p3)--(p2) (p4)--(p2);
  \node[lb,anchor=west,align=left] at (-0.1,-0.10) {papers only; no authors};
\end{scope}
% ---- (b) co-authorship
\begin{scope}[xshift=5.0cm]
  \node[tb,anchor=west] at (-0.1,2.15) {(b) co-authorship};
  \node[atm] (a1) at (0.35,1.40) {}; \node[atm] (a2) at (1.45,1.55) {};
  \node[atm] (a3) at (0.90,0.70) {}; \node[atm] (a4) at (2.15,0.90) {};
  \draw[inc] (a1)--(a2)--(a3)--(a1) (a2)--(a4);
  \node[lb,anchor=west,align=left] at (-0.1,-0.10) {authors only; the paper is gone};
\end{scope}
% ---- (c) bipartite
\begin{scope}[yshift=-3.1cm]
  \node[tb,anchor=west] at (-0.1,2.15) {(c) bipartite};
  \node[cpx] (q1) at (0.55,1.55) {}; \node[cpx] (q2) at (1.85,1.55) {};
  \node[atm] (b1) at (0.25,0.65) {}; \node[atm] (b2) at (1.20,0.65) {};
  \node[atm] (b3) at (2.15,0.65) {};
  \draw[inc] (q1)--(b1) (q1)--(b2) (q2)--(b2) (q2)--(b3);
  \node[lb,anchor=west,align=left] at (-0.1,-0.10) {both kinds of thing; no citations};
\end{scope}
% ---- (d) author hypergraph
\begin{scope}[xshift=5.0cm,yshift=-3.1cm]
  \node[tb,anchor=west] at (-0.1,2.15) {(d) author hypergraph};
  \node[nd] (c1) at (0.40,1.55) {}; \node[nd] (c2) at (1.35,1.60) {};
  \node[nd] (c3) at (0.85,0.75) {}; \node[nd] (c4) at (2.25,1.05) {};
  \begin{scope}[on background layer]
    \node[cx,fill=black!26,fill opacity=0.7,fit=(c1)(c2)(c3)]{};
    \node[cx,fill=black!26,fill opacity=0.7,inner sep=3.6pt,fit=(c2)(c4)]{};
  \end{scope}
  \node[lb,anchor=west,align=left] at (-0.1,-0.10) {the paper is one object; no citations};
\end{scope}
\end{tikzpicture}
\end{adjustbox}
\caption{\textbf{Four maps of one system.} The scientific literature drawn four
ways. Filled circles are authors, rectangles and shaded boxes are papers, solid
lines are authorship and dashed lines citation. Each map is made by discarding
something that was in the system: (a) forgets who wrote anything, (b) and (d)
forget citation, and (b) also forgets that a paper with five authors is one act
and not ten. No two of them are the same object, so nothing proved on one can be
quoted on another. Where two complexes overlap, as in (d), the shading
compounds, so a darker patch marks the atoms they share --- a convention used
throughout the book.}
\label{fig:projections}
\end{figure}

A single paper is neither a node nor a link. It is a small structured thing in its own right: it has an author list,
and it has a reference list. The author list is a set of people. The reference
list is a set of \emph{other papers}. And the two lists do not mix --- no author
appears in the bibliography as a reference, and no reference is a co-author.

Suppose we agree to call the paper one object and its two lists the internal structure of that object.
Then the paper's interior is not a single group but two disjoint groups: a
person entering the paper through the author list can reach the other authors
and stops there, and a paper entering through the reference list can reach the
other references and stops there. To get from an author of a paper to something
that paper cites, one has to leave the paper and re-enter by the other door.
Whether or not that is how citation actually works sociologically, it is
certainly how the document is built, and any representation that flattens the
paper into one undifferentiated blob of ``things associated with this paper''
has lost it.

So the literature demands two things at once, and neither of the standard
languages provides both. It demands that a complex object be able to
\emph{contain another complex object of the same kind} --- a paper contains
papers, in its reference list. And it demands that the interior of that object
have \emph{structure of its own} --- not one list but two, and the difference
between them matters. Everything in this chapter is an attempt to meet those
two demands and nothing more.

\begin{figure}[t]
\centering
\begin{adjustbox}{max width=\linewidth}
\begin{tikzpicture}[yscale=1.15]
% three publications, each a complex with two internal hyperedges
\foreach \x/\n in {0/1, 3.1/2, 6.2/3} {
  \node[cpx] (A\n) at (\x,0) {};
  \node[cpxb] (R\n) at (\x+1.0,0) {};
  \begin{scope}[on background layer]
    \node[pub,fit=(A\n)(R\n)] (P\n) {};
  \end{scope}
}
\node[lb,anchor=north east] at (-0.18,-0.10) {$A_i$};
\node[lb,anchor=north west] at (1.18,-0.10) {$R_i$};
% authors
\foreach \x/\m in {-0.35/1, 0.75/2, 1.95/3, 3.15/4, 4.35/5, 5.55/6, 6.75/7, 7.85/8}
  \node[atm] (u\m) at (\x,-1.75) {};
\draw[inc] (A1)--(u1) (A1)--(u2) (A1)--(u3) (A2)--(u3) (A2)--(u4) (A2)--(u5)
           (A3)--(u6) (A3)--(u7) (A3)--(u8);
% citations: the cited publication is a vertex of the citing publication's reference list
\draw[incd] (P2) to[out=125,in=55] (R1);
\draw[incd] (P3) to[out=125,in=55] (R2);
\draw[incd] (P3) to[out=140,in=48] (R1);
% legend, laid out below rather than beside so the drawing keeps the page width
\node[cpx]  (l1) at (0.15,-2.95) {}; \node[lb,anchor=west] at (0.42,-2.95) {author list $A_i$};
\node[cpxb] (l2) at (3.25,-2.95) {}; \node[lb,anchor=west] at (3.52,-2.95) {reference list $R_i$};
\node[atm]  (l3) at (6.55,-2.95) {}; \node[lb,anchor=west] at (6.82,-2.95) {author (atom)};
\draw[inc]  (0.05,-3.50)--(0.55,-3.50); \node[lb,anchor=west] at (0.72,-3.50) {authorship};
\draw[incd] (3.15,-3.50)--(3.65,-3.50); \node[lb,anchor=west] at (3.82,-3.50) {citation};
\draw[rounded corners=1.6pt,draw=black!70,line width=.6pt]
  (6.45,-3.62) rectangle (7.05,-3.38);
\node[lb,anchor=west] at (7.22,-3.50) {publication (complex)};
\end{tikzpicture}
\end{adjustbox}
\caption{\textbf{The literature as a chygraph} \citep{vazquez2023pre}. A
publication is a complex, drawn as a rounded box. Its interior is a hypergraph
with two hyperedges: the author list $A_i$ (white) and the reference list $R_i$
(shaded). Authors are atoms and are included in author lists; a citation is the
inclusion of one publication in another publication's reference list, so a
complex is a vertex of a complex. The two hyperedges do not overlap, so the
interior has two components, and the door of entry decides what is reachable.}
\label{fig:publications}
\end{figure}

\section{Graphs, hypergraphs and ubergraphs}

A \emph{graph} $G(V,E)$ is a set of vertices $V$ and a set of edges $E$, each
edge a pair of vertices. It meets neither demand. An edge cannot contain an
edge, and an edge has no interior to give structure to --- it is a pair, and a
pair is exhausted by naming its two members.

A \emph{hypergraph} $H(V,E)$ relaxes the pair: a hyperedge is any subset of $V$,
so an interaction among five things is one hyperedge rather than ten links
\citep{ghoshal2009,battiston2020}. That is real progress, and it is the first
move of the modern higher-order-network literature. It meets the second
demand only barely --- a hyperedge is a flat set, with no internal organisation
--- and it does not meet the first at all. Hyperedges contain vertices. They do
not contain hyperedges. Our paper cannot cite.

The first demand was met by Joslyn and Nowak, who defined \emph{ubergraphs}
\citep{joslyn2017}. We restate their construction in the language this book uses
throughout. Begin by dividing the things in a system into
those we have chosen not to decompose any further --- call them \emph{atoms} ---
and those that are made of other things --- call them \emph{complexes}. Which is
which is a modelling decision, not a fact about the world: an author is an atom
here because we have decided not to model people as made of parts, and nothing
stops someone else deciding otherwise. Then:

\begin{quote}
\emph{An ubergraph} $\mathcal{U}(A,C)$ is a set of atoms $A$ and a set of
complexes $C$, where each complex is a subset of $A\cup C$.
\end{quote}

The clause that does the work is $A\cup C$ rather than $A$. A complex may
contain complexes. Our paper can now cite --- and it still cannot tell its
authors from its references, because a complex is still a flat set. Both lists
are merged into one.

\section{Chygraph definition}
\label{sec:definition}
\index{chygraph|textbf}
\index{complex|textbf}
\index{atom|textbf}
\index{inclusion}

The remaining step is small. If the trouble
is that a complex's interior is a flat set, give the interior the same treatment
we have just given the exterior: make it a hypergraph.

\begin{quote}
\emph{A complex hypergraph} (chygraph) $\chi(A,C)$ is a set of atoms $A$ and a
set of complexes $C$, where each complex is a \emph{hypergraph} whose vertex set
lies in $A\cup C$ \citep{vazquez2023pre}.
\end{quote}

There is a tighter way to say it that will be more convenient from
Chapter~\ref{ch:percolation} onward. An atom is nothing but a complex that
happens to include no complexes, so we need not carry $A$ separately:

\begin{quote}
A chygraph is a set of complexes, where complexes are hypergraphs with a vertex
set in the set of complexes \citep{vazquez2024comnet}.
\end{quote}

Two properties are packed into that sentence. \emph{Self-reference:} complexes
are built from complexes, and are themselves available as building blocks.
\emph{Fine-grained structure:} what is inside a complex is not a heap but a
hypergraph.

Figure~\ref{fig:ladder} puts the three constructions side by side. Each step
buys exactly one thing, and the book needs both.

\begin{figure}[t]
\centering
\begin{adjustbox}{max width=\linewidth}
\begin{tikzpicture}[yscale=1.0]
\begin{scope}
  \node[tb,anchor=west] at (-0.2,2.65) {hypergraph};
  \node[cpx] (h) at (0.9,1.25) {};
  \node[atm] (h1) at (0.1,0.45) {}; \node[atm] (h2) at (0.9,0.45) {};
  \node[atm] (h3) at (1.7,0.45) {};
  \draw[inc] (h)--(h1) (h)--(h2) (h)--(h3);
  \node[lb,anchor=west,align=left] at (-0.2,-0.35)
    {a complex holds atoms;\\ its interior is a flat set};
\end{scope}
\begin{scope}[xshift=4.35cm]
  \node[tb,anchor=west] at (-0.2,2.65) {ubergraph};
  \node[cpx] (u) at (0.9,1.95) {};
  \node[cpx] (v) at (0.15,1.25) {};
  \node[atm] (u1) at (1.75,1.25) {};
  \node[atm] (v1) at (-0.3,0.45) {}; \node[atm] (v2) at (0.6,0.45) {};
  \draw[inc] (u)--(v) (u)--(u1) (v)--(v1) (v)--(v2);
  \node[lb,anchor=west,align=left] at (-0.2,-0.35)
    {a complex may hold complexes;\\ its interior is still flat};
\end{scope}
\begin{scope}[xshift=8.9cm]
  \node[tb,anchor=west] at (-0.2,2.65) {chygraph};
  \node[cpx]  (e1) at (0.55,1.95) {};
  \node[cpxb] (e2) at (1.35,1.95) {};
  \begin{scope}[on background layer]\node[pub,fit=(e1)(e2)] (W) {};\end{scope}
  \node[cpx] (w) at (2.05,1.20) {};
  \node[atm] (w1) at (0.05,1.20) {}; \node[atm] (w2) at (0.95,1.20) {};
  \draw[inc] (e1)--(w1) (e1)--(w2) (e2)--(w);
  \node[lb,anchor=west,align=left] at (-0.2,-0.35)
    {a complex holds complexes,\\ \emph{and} its interior is a hypergraph};
\end{scope}
\end{tikzpicture}
\end{adjustbox}
\caption{\textbf{Three constructions, two demands.} A hypergraph's complexes
hold atoms only. An ubergraph's complexes may hold other complexes, but a
complex is still a flat set. A chygraph's complexes hold complexes \emph{and}
carry a hypergraph inside, so a publication can distinguish its author list
(white) from its reference list (shaded). An ubergraph is the special case of a
chygraph in which every interior hypergraph has exactly one hyperedge, so every
result in this book applies to ubergraphs with ``component size'' read as
``cardinality''.}
\label{fig:ladder}
\end{figure}

The literature is now representable without loss. Writing it out,
\begin{equation}
  \chi\Big(A,\ \big\{\,\mathcal{H}_i\big(A_i\cup R_i,\ \{A_i,R_i\}\big)\,\big\}\Big),
  \label{eq:publications}
\end{equation}
where $A$ is the set of authors, $i$ runs over publications, $A_i$ is the
author list of publication $i$ and $R_i$ its reference list. Each publication is
a hypergraph on the vertex set $A_i\cup R_i$ with exactly two hyperedges,
$A_i$ and $R_i$ (Fig.~\ref{fig:publications}). Nothing has been discarded: the
citation network of Fig.~\ref{fig:projections}a, the co-authorship network of
Fig.~\ref{fig:projections}b, the bipartite graph of
Fig.~\ref{fig:projections}c and the hypergraph of Fig.~\ref{fig:projections}d
are all recoverable from Eq.~\eqref{eq:publications} by forgetting, and none of
them can recover any of the others.

\section{The vocabulary}
\label{sec:vocabulary}
\index{cardinality|textbf}
\index{chy-degree|textbf}
\index{excess cardinality|textbf}

A handful of words, all of them intended to be as flat-footed as they sound.

Complex $A$ \emph{includes} complex $B$ when $B$ is in the vertex set of $A$;
equivalently, $B$ is \emph{included in} $A$. An \emph{atom} is a complex that
includes no complexes. The \emph{cardinality} of a complex is the number of
complexes it includes. The \emph{chy-degree} of a complex is the number of
complexes that include it.

Two conventions. First, quantities defined at the
chygraph level get Greek letters, to keep them separate from the same quantities
computed on some graph made out of the chygraph; so the chy-degree is $\kappa$
and never $k$. Second, ``cardinality'' counts what is inside and ``chy-degree''
counts what is outside, and confusing them is the single commonest slip in this
subject. In the literature example, an author's chy-degree is the number of
papers they have written, and a publication's chy-degree is the number of
citations it has received; a publication's cardinality is its number of authors
plus its number of references. An author's $h$-index is a statistic of the
chy-degrees of the complexes that include that atom.

Formally, index the atoms and complexes together as $i=1,\dots,n$ with
$n=|A\cup C|$, and let $C_j$ denote the vertex set of complex $j$. The
\emph{chy-adjacency matrix} is
\begin{equation}
  \alpha_{ij}=\begin{cases}1 & \text{if } i\in C_j,\\ 0 & \text{otherwise,}\end{cases}
  \label{eq:alpha}
\end{equation}
and the chy-degree and cardinality are its two marginals,
\begin{equation}
  \kappa_i=\sum_j\alpha_{ij},\qquad c_j=\sum_i\alpha_{ij}.
  \label{eq:kappa}
\end{equation}
That $\alpha$ is not symmetric --- unlike the adjacency matrix of a graph, which
it otherwise resembles --- is the difference between being inside something and
containing something.

\subsection{The two ways out of a complex}

The rest of the book rests on one sentence, which every recursion in
Chapters~\ref{ch:percolation} to~\ref{ch:overlap} formalises.

\begin{quote}
Standing at a complex, there are exactly two ways to leave it: \emph{up},
through a complex that includes it, or \emph{down}, into the complexes it
includes.
\end{quote}

Up, there are $\kappa$ choices --- one per inclusion --- and they lead to
distinct complexes elsewhere in the chygraph. Down, there is the interior
hypergraph, and this is where the second of the two demands matters:
going down is not a matter of choosing one of $c$ members at random, because the
interior may be several disjoint pieces and only the piece entered is
reachable. Entering a publication as one of its authors reaches the other
authors; entering it as one of its references reaches the other references.
The relevant number is therefore not the cardinality but the size of the
\emph{intra-complex component} arrived in, written $s$
(Fig.~\ref{fig:twoways}).

\begin{figure}[t]
\centering
\begin{adjustbox}{max width=0.92\linewidth}
\begin{tikzpicture}
% the complex in the middle
\node[cpx] (E1) at (2.0,0) {}; \node[cpxb] (E2) at (3.1,0) {};
\begin{scope}[on background layer]\node[pub,fit=(E1)(E2)] (P) {};\end{scope}
\node[lb,anchor=west] at (3.55,0.0) {a complex};
% up: complexes that include it
\node[cpx] (Ua) at (1.15,1.5) {}; \node[cpx] (Ub) at (2.55,1.75) {};
\node[cpx] (Uc) at (3.95,1.5) {};
\draw[msg] (2.30,0.34) -- (1.35,1.30);
\draw[msg] (2.60,0.36) -- (2.57,1.55);
\draw[msg] (2.90,0.34) -- (3.75,1.30);
\node[lb,anchor=south,align=center] at (2.55,2.05)
  {\emph{up:} $\kappa$ complexes include it};
% down: the interior, two components
\node[atm] (d1) at (1.15,-1.6) {}; \node[atm] (d2) at (2.0,-1.6) {};
\node[atm] (d3) at (2.85,-1.6) {};
\node[cpx] (d4) at (3.85,-1.6) {}; \node[cpx] (d5) at (4.75,-1.6) {};
\draw[inc] (E1)--(d1) (E1)--(d2) (E1)--(d3) (E2)--(d4) (E2)--(d5);
\draw[decorate,decoration={brace,amplitude=3pt,mirror},black!55]
  (1.0,-1.9) -- (3.0,-1.9);
\draw[decorate,decoration={brace,amplitude=3pt,mirror},black!55]
  (3.7,-1.9) -- (4.9,-1.9);
\node[lb,anchor=north] at (2.0,-2.05) {component, $s=3$};
\node[lb,anchor=north] at (4.3,-2.05) {component, $s=2$};
\node[lb,anchor=west,align=left] at (5.35,-1.30)
  {\emph{down:} into the interior ---\\ but only the component\\ entered by};
\end{tikzpicture}
\end{adjustbox}
\caption{\textbf{The two ways out.} Up, through any of the $\kappa$ complexes
that include this one. Down, into its interior hypergraph --- where what is
reachable is not the whole cardinality $c=5$ but the size $s$ of the component
entered. A publication's author list and reference list are two components of
size $3$ and $2$; arriving as an author, the two other authors are reachable,
not the four other members. Percolation (Chapter~\ref{ch:percolation}) is this
picture with a probability attached to each arrow; the cavity recursion
(Chapter~\ref{ch:statmech}) is this picture with a message.}
\label{fig:twoways}
\end{figure}

Cardinality and component size coincide whenever a complex's interior hypergraph
has a single hyperedge --- which is to say, on ubergraphs, and therefore on
graphs, hypergraphs and every structure in Sec.~\ref{sec:everything} below.
They part company only for genuinely compound complexes like our publications.
The theory is written in terms of $s$ throughout, so both cases are covered by
one calculation, and a reader who only ever cares about hypergraphs may read $s$
as the cardinality, and its excess $\bar s$ of Eq.~\eqref{eq:excess} as
``cardinality minus the one arrived by'', and lose nothing.

\section{Layers}
\label{sec:layers}
\index{layer|textbf}

One more piece of bookkeeping is needed, and it is what makes the construction
computationally workable.

The complexes of a chygraph are usually not statistically alike. In the
literature example, publications and authors are different sorts of object with
different chy-degree distributions, and inside a publication the author list and
the reference list behave differently too. In a multiplex network, a friendship
link and a co-worker link join the same people and are not interchangeable. So
we partition the complexes into \emph{layers}, indexed $l=0,\dots,L-1$, with the
understanding that complexes within a layer are statistically alike and
complexes in different layers need not be.

That is the working definition, and for every calculation in this book it is
all that is needed. The formal version, which fixes how deep a chygraph is
rather than merely how many kinds of thing it contains, is due to
\citet{vazquez2023pre}:

\begin{quote}
\emph{Chygraph length.} Let $\Pi=\Pi_1\cup\dots\cup\Pi_l$ be a partition of the
complex set that is non-intersecting ($\Pi_i\cap\Pi_j=\emptyset$ for $i\ne j$),
hierarchical (if a complex lies in $\Pi_j$ then its vertices lie in
$A\cup(\cup_{k\le j}\Pi_k)$), and statistically homogeneous within each part.
The chygraph length $L(\chi)$ is the largest $l$ over all such partitions.
\end{quote}

The hierarchical clause says that a layer may only be built out of layers below
it, which is what makes ``levels of inclusion'' well defined. Human geography is
the everyday example: people, neighbourhoods, cities, countries, continents ---
five levels, each assembled from the one beneath. The homogeneity clause is what
lets two layers sit at the same level and still be told apart, which is what a
multiplex network needs. The two clauses pull in opposite directions --- one
wants few, deep layers, the other wants many, distinguishable ones --- and the
maximum in the definition resolves them.

No calculation in this book requires computing $L(\chi)$ for a given chygraph. In practice one \emph{chooses} the layers when
setting up the model, and the choice is part of the modelling, exactly as
choosing what counts as an atom was.

\section{Everything is a chygraph}
\label{sec:everything}
\index{hypergraph}
\index{multiplex network}
\index{graph, as a chygraph}

Each structure below is obtained by writing down its layers and their
cardinalities, and nothing else changes. Figure~\ref{fig:mappings} draws them in
a single vocabulary.

\paragraph{A graph} is two layers. Layer $0$ is atoms, one per node. Layer $1$
is complexes of cardinality two, one per link, each containing its two
endpoints. A node's chy-degree is its ordinary degree; a link's chy-degree is
zero, because nothing includes a link.

\paragraph{A hypergraph} is the same with the cardinality freed: layer $1$
complexes contain any number of atoms. A uniform hypergraph fixes that number at
$c$; a clique network --- a graph built by pasting $c$-cliques together --- is
the same chygraph again, because a clique and a hyperedge differ in how they are
drawn and not in what includes what.

\paragraph{A multiplex network} is one atom layer and one complex layer per link
type: friendship in layer $1$, co-working in layer $2$, and so on
\citep{gomez2013,kivela2014}. The atoms are shared, which is exactly what makes
it a multiplex rather than two unrelated networks, and no further apparatus is
required.

\paragraph{Interacting graphs and hypergraphs} need one more layer than one
might guess. Two hypergraphs that interact are: an atom layer and a hyperedge
layer for each, \emph{plus} a layer of hyperedges containing atoms from both
\citep{leicht2009,bianconi2017}. With more than two, one adds a cross layer per
interacting pair.

\paragraph{A clustered graph} becomes a chygraph by \emph{promoting its dense
motifs to complexes}. This is the mapping Part~III rests on. A graph rich in triangles is given two complex layers, one of
cardinality two for the plain links and one of cardinality three for the
triangles \citep{newman2009,karrer2010}. A graph rich in maximal cliques of
assorted sizes gets one layer per size. The triangle has stopped being three
links that happen to close and has become one object, and the loop that made it
a triangle has gone inside the object where the theory can sum over it exactly.
That move is what Chapter~\ref{ch:statmech} is built on, and its price is stated
in Sec.~\ref{sec:treelike} below.

\begin{figure}[t]
\centering
\begin{adjustbox}{max width=\linewidth}
\begin{tikzpicture}
% ---------- row 1: graph, hypergraph, multiplex ----------
\begin{scope}
  \node[tb,anchor=west] at (-0.3,1.75) {graph};
  \node[cpx] (g1) at (0.35,1.05) {}; \node[cpx] (g2) at (1.45,1.05) {};
  \node[atm] (n1) at (-0.05,0.25) {}; \node[atm] (n2) at (0.85,0.25) {};
  \node[atm] (n3) at (1.85,0.25) {};
  \draw[inc] (g1)--(n1) (g1)--(n2) (g2)--(n2) (g2)--(n3);
  \node[lb,anchor=west] at (-0.3,-0.35) {$L=2$, $c=2$};
\end{scope}
\begin{scope}[xshift=3.6cm]
  \node[tb,anchor=west] at (-0.3,1.75) {hypergraph};
  \node[cpx] (H1) at (0.55,1.05) {}; \node[cpx] (H2) at (1.85,1.05) {};
  \node[atm] (m1) at (-0.15,0.25) {}; \node[atm] (m2) at (0.55,0.25) {};
  \node[atm] (m3) at (1.25,0.25) {}; \node[atm] (m4) at (2.25,0.25) {};
  \draw[inc] (H1)--(m1) (H1)--(m2) (H1)--(m3) (H2)--(m3) (H2)--(m4);
  \node[lb,anchor=west] at (-0.3,-0.35) {$L=2$, $c$ free};
\end{scope}
\begin{scope}[xshift=7.5cm]
  \node[tb,anchor=west] at (-0.3,1.75) {multiplex};
  \node[cpx]  (M1) at (0.35,1.20) {}; \node[cpx]  (M2) at (1.65,1.20) {};
  \node[atm] (p1) at (-0.05,0.45) {}; \node[atm] (p2) at (0.90,0.45) {};
  \node[atm] (p3) at (2.05,0.45) {};
  \node[cpxc] (M3) at (0.45,-0.30) {}; \node[cpxc] (M4) at (1.75,-0.30) {};
  \draw[inc] (M1)--(p1) (M1)--(p2) (M2)--(p2) (M2)--(p3)
             (M3)--(p1) (M3)--(p3) (M4)--(p2) (M4)--(p3);
  \node[lb,anchor=west] at (-0.3,-0.90) {one layer per link type};
\end{scope}
% ---------- row 2: interacting, clustered graph ----------
\begin{scope}[yshift=-3.5cm,xshift=0.6cm]
  \node[tb,anchor=west] at (-0.3,1.75) {interacting};
  \node[cpx]  (I1) at (0.30,1.05) {}; \node[cpxb] (I3) at (1.35,1.05) {};
  \node[cpxc] (I2) at (2.35,1.05) {};
  \node[atm] (q1) at (-0.05,0.25) {}; \node[atm] (q2) at (0.80,0.25) {};
  \node[nd]  (q3) at (1.85,0.25) {}; \node[nd]  (q4) at (2.70,0.25) {};
  \draw[inc] (I1)--(q1) (I1)--(q2) (I2)--(q3) (I2)--(q4)
             (I3)--(q2) (I3)--(q3);
  \node[lb,anchor=west,align=left] at (-0.3,-0.48)
    {a cross layer (mid grey) per\\ interacting pair};
\end{scope}
\begin{scope}[yshift=-3.5cm,xshift=5.9cm]
  \node[tb,anchor=west] at (-0.3,1.75) {clustered graph};
  \node[cpx]  (T1) at (0.35,1.05) {};
  \node[cpxb] (T2) at (1.85,1.05) {};
  \node[atm] (r1) at (-0.15,0.25) {}; \node[atm] (r2) at (0.85,0.25) {};
  \node[atm] (r3) at (1.75,0.25) {}; \node[atm] (r4) at (2.65,0.25) {};
  \draw[inc] (T1)--(r1) (T1)--(r2) (T2)--(r2) (T2)--(r3) (T2)--(r4);
  \node[lb,anchor=west,align=left] at (-0.3,-0.48)
    {links at $c=2$, triangles at\\ $c=3$: motifs promoted};
\end{scope}
\end{tikzpicture}
\end{adjustbox}
\caption{\textbf{One object, five structures.} Circles are atoms, rectangles
are complexes, and shading distinguishes layers; a line is an inclusion. In the
interacting panel the two hypergraphs' atoms are told apart by fill, and the
mid-grey complexes are the cross layer that contains atoms of both. Nothing changes from panel to panel except which layers exist and
what cardinalities they carry. This is the claim the rest of the book relies on:
a calculation done once on the chygraph is a calculation done for all five, and
the specialisation is substitution rather than rederivation.}
\label{fig:mappings}
\end{figure}

The list is the argument for the whole construction. The
higher-order-network literature has produced a separate percolation calculation
for hypergraphs, for multiplex networks, for interacting networks, for
interdependent networks, for simplicial complexes, and for clustered graphs, and
the number of independent derivations has been growing faster than anyone's
ability to hold them in one head \citep{battiston2020,bianconi2021}. If all six
are chygraphs, one calculation covers them, and the differences between them
become differences in what one substitutes at the end. Chapter~\ref{ch:percolation}
carries that out.

\section{Two ways to draw the same thing}
\label{sec:drawing}
\index{inclusion diagram}

Figures~\ref{fig:publications} and~\ref{fig:mappings} draw a complex as a
\emph{node} of its own, with lines running to what it includes. Every figure
after this chapter draws a complex as a \emph{box around} what it includes.
These are the same object. Figure~\ref{fig:twodrawings} puts the two side by
side once, so that neither reads as a new construction later.

The inclusion drawing is complete and scales badly: it puts every complex and
every inclusion on the page, and it becomes unreadable at a couple of dozen
complexes. The spatial drawing is compact and incomplete: it can only show a
complex as a box when the complex's members can be drawn adjacent, which
fails as soon as two complexes share more than one atom --- precisely the
situation of Sec.~\ref{sec:treelike}. Use the first when the question is what
includes what, the second when the question is what is near what.

\begin{figure}[t]
\centering
\begin{adjustbox}{max width=0.9\linewidth}
\begin{tikzpicture}
\begin{scope}
  \node[tb,anchor=west] at (-0.2,2.15) {(a) inclusion drawing};
  \node[cpx] (A) at (0.6,1.45) {}; \node[cpx] (B) at (2.1,1.45) {};
  \node[lb,anchor=south] at (0.6,1.62) {$a$};
  \node[lb,anchor=south] at (2.1,1.62) {$b$};
  \node[atm] (x1) at (0.0,0.55) {}; \node[atm] (x2) at (1.2,0.55) {};
  \node[atm] (x3) at (2.7,0.55) {};
  \draw[inc] (A)--(x1) (A)--(x2) (B)--(x2) (B)--(x3);
  \node[lb,anchor=west,align=left] at (-0.2,-0.20)
    {complexes are nodes;\\ lines are inclusions};
\end{scope}
\begin{scope}[xshift=5.6cm]
  \node[tb,anchor=west] at (-0.2,2.15) {(b) spatial drawing};
  % the two complexes are drawn at different heights and different fills, so
  % that the atom they share is visibly in both rather than in one long blob
  \node[nd] (y1) at (0.15,1.05) {}; \node[nd] (y2) at (1.35,1.45) {};
  \node[nd] (y3) at (2.55,1.05) {};
  \begin{scope}[on background layer]
    \node[cx,fill=black!26,fill opacity=0.7,fit=(y1)(y2)] {};
    \node[cx,fill=black!26,fill opacity=0.7,inner sep=3.8pt,fit=(y2)(y3)] {};
  \end{scope}
  \node[lb,anchor=east] at (0.05,1.62) {$a$};
  \node[lb,anchor=west] at (2.70,1.62) {$b$};
  \node[lb,anchor=west,align=left] at (-0.2,-0.20)
    {complexes are boxes\\ drawn around their members};
\end{scope}
\end{tikzpicture}
\end{adjustbox}
\caption{\textbf{The same chygraph twice.} Three atoms and two complexes, $a$
and $b$, the middle atom belonging to both. The inclusion drawing (a) is used in
this chapter because the question here is what includes what; the spatial
drawing (b) is used everywhere afterwards because the question there is what a
message can reach. Neither adds anything to the other.}
\label{fig:twodrawings}
\end{figure}

\section{Random chygraphs}
\label{sec:ensembles}
\index{random chygraph|textbf}
\index{ensemble}

The rest of the book computes averages over ensembles rather than properties of
one object, so we need to say what specifying an ensemble amounts to. It amounts
to specifying, for every ordered pair of layers, how many complexes of one layer
attach to a complex of the other, in each of the two directions of
Sec.~\ref{sec:vocabulary} --- and that is all.

Let $\Phi^l_k(x)$ be the probability generating function of the number of
layer-$k$ complexes that include a layer-$l$ complex, and $G^l_k(y)$ that of the
number of layer-$k$ complexes in the intra-complex component of a layer-$l$
complex, entered at a given vertex. Write $\bar\Phi^l_k$ and $\bar G^l_k$ for
the corresponding \emph{excess} functions --- the same counts taken from a
complex reached along an inclusion, with the inclusion used to arrive not
counted again. Up and down, ordinary and excess: four families of functions, one
per direction per situation.

Their first derivatives at $1$ are the four moment matrices that will run
through every chapter:
\begin{align}
  \Phi^{l\prime}_k(1)&=\ave{\kappa}_{lk}, &
  \bar\Phi^{l\prime}_k(1)&=\ave{\kbar}_{lk},\nonumber\\
  G^{l\prime}_k(1)&=\ave{s}_{lk}, &
  \bar G^{l\prime}_k(1)&=\ave{\sbar}_{lk},
  \label{eq:moments}
\end{align}
with the excess averages given, as usual for a size-biased count, by
\begin{equation}
  \ave{\kbar}_{lk}=\frac{\ave{\kappa(\kappa-1)}_{lk}}{\ave{\kappa}_{lk}},
  \qquad
  \ave{\sbar}_{lk}=\frac{\ave{s(s-1)}_{lk}}{\ave{s}_{lk}}.
  \label{eq:excess}
\end{equation}

Four $L\times L$ matrices, then, summarise a chygraph at the level of first
moments: $\ave{\kappa}$, $\ave{\kbar}$, $\ave{s}$, $\ave{\sbar}$. Whatever the sizes of the complexes, whatever
the depth of the inclusion hierarchy, the first-moment description of the
ensemble is four matrices whose dimension is the number of \emph{layers} --- two
for a graph, three for a two-layer multiplex --- and not the number of complexes
or their cardinality. Chapter~\ref{ch:percolation} shows that the percolation
threshold of every structure in Sec.~\ref{sec:everything} is fixed by these four
matrices alone; Chapter~\ref{ch:giant} shows precisely what they fail to
determine, which is everything about the transition except where it is.

Two caveats, stated here so they are not sprung later. Equation~\eqref{eq:moments}
takes the generating functions to factorise over target layers --- a complex's
participation in one layer independent of its participation in another --- which
covers every mapping in Sec.~\ref{sec:everything} but is a genuine restriction,
lifted in Sec.~\ref{sec:joint}. And moments are moments: two ensembles can agree
on all four matrices and disagree about the shape of the transition, which is
the subject of Chapter~\ref{ch:giant} and not a defect to be repaired.

\section{Treelike over complexes, and what it costs}
\label{sec:treelike}
\index{treelike|textbf}
\index{overlap!between complexes}

We close the chapter with the property that motivates everything after
Chapter~\ref{ch:percolation}, and with its price, which is structural rather
than a matter of care.

The methods of this book --- generating functions for percolation, the cavity
recursion for statistical mechanics --- require that the object they run on be
locally treelike: that following inclusions away from a complex, one does not
come back. Real graphs are not treelike: the neighbours of a node are neighbours
of one another, and the triangle they form is a loop of length three, the
shortest there is.

Promote that triangle to a complex, though, and the loop is gone from the
structure \emph{over} complexes. It has not been discarded --- it is inside the
complex, where Chapter~\ref{ch:statmech} sums over it exactly at a cost that
depends on the cardinality and on nothing else. What was an obstruction has
become an interior. This is the bargain the book explores, and Fig.~\ref{fig:price}a
is the good case.

Figure~\ref{fig:price}b is the bad one. If two complexes share \emph{two} atoms
rather than one, then there is a loop that promotion did not absorb: leave the
first complex through one shared atom, cross the second complex, and come back
through the other. Over complexes, that is a loop of length two, and the
recursion will double-count it.

\paragraph{The pairwise test is necessary and not sufficient.} It is convenient
to read the condition as the single question ``do two complexes ever share two
atoms?'', and that is how the rest of this book applies it. On an ensemble the
reading is safe: the random chygraphs of Part~II are locally treelike in the
full sense, and a loop of any length is a rare event. On a \emph{given} network
it is not. A loop can run through three or more complexes, each meeting the next
in a single atom and no pair sharing two, so that the pairwise question is
answered ``no'' and the structure over complexes still closes on itself. Four
triangles glued in a ring at single vertices do exactly that, and the recursion
is wrong on them; Section~\ref{sec:mergeerror} draws the object and prices the
error. What the recursion needs is that the incidence structure be a forest, of
which sharing at most one atom is a consequence rather than a guarantee.

\begin{figure}[t]
\centering
\begin{adjustbox}{max width=0.94\linewidth}
\begin{tikzpicture}
\begin{scope}
  \node[tb,anchor=west] at (-0.2,2.55) {(a) meeting in one atom};
  \node[nd] (a1) at (0.15,1.35) {}; \node[nd] (a2) at (1.10,1.85) {};
  \node[nd] (a3) at (1.00,0.80) {};
  \node[nd] (a4) at (2.30,1.95) {}; \node[nd] (a5) at (2.45,1.05) {};
  \draw[lnk] (a1)--(a2)--(a3)--(a1) (a2)--(a4)--(a5)--(a2);
  \begin{scope}[on background layer]
    \node[cx,fill=black!26,fill opacity=0.7,fit=(a1)(a2)(a3)]{};
    \node[cx,fill=black!26,fill opacity=0.7,inner sep=3.8pt,fit=(a2)(a4)(a5)]{};
  \end{scope}
  \node[lb,anchor=east] at (0.02,1.90) {$a$};
  \node[lb,anchor=west] at (2.72,1.95) {$b$};
  \node[lb,anchor=west,align=left] at (-0.2,0.05)
    {treelike over complexes:\\ no loop survives promotion};
\end{scope}
\begin{scope}[xshift=5.5cm]
  \node[tb,anchor=west] at (-0.2,2.55) {(b) sharing two atoms};
  \node[nd] (b1) at (0.15,1.35) {}; \node[nd] (b2) at (1.20,1.95) {};
  \node[nd] (b3) at (1.20,0.80) {}; \node[nd] (b4) at (2.30,1.35) {};
  \draw[lnk] (b1)--(b2)--(b3)--(b1) (b2)--(b4)--(b3);
  \begin{scope}[on background layer]
    \node[cx,fill=black!26,fill opacity=0.7,fit=(b1)(b2)(b3)]{};
    \node[cx,fill=black!26,fill opacity=0.7,inner sep=3.8pt,fit=(b2)(b3)(b4)]{};
  \end{scope}
  \node[lb,anchor=east] at (0.02,1.95) {$a$};
  \node[lb,anchor=west] at (2.62,1.35) {$b$};
  % the surviving loop: out of a through one shared atom, back through the other
  \draw[msg,black!85] (1.05,1.72) to[out=250,in=110] (1.05,1.03);
  \draw[msg,black!85] (1.35,1.03) to[out=70,in=290] (1.35,1.72);
  \node[lb,anchor=west,align=left] at (-0.2,0.05)
    {a loop of length two survives:\\ the recursion double-counts};
\end{scope}
\end{tikzpicture}
\end{adjustbox}
\caption{\textbf{The bargain and its price.} (a) Two triangles, $a$ and $b$,
meeting at one atom. Promoted to complexes, the structure over complexes is a
tree, and every loop of the underlying graph now lives inside a complex.
(b) Two triangles sharing an edge --- two atoms. Promotion removes the loops
inside each triangle and leaves a loop \emph{between} them, traced by the
arrows, which the theory of Chapters~\ref{ch:percolation}--\ref{ch:satisfiability}
does not see. Chapter~\ref{ch:overlap} measures what this costs on real
clustered graphs, and Chapters~\ref{ch:gbp} and~\ref{ch:metacomplex} are the two
repairs.}
\label{fig:price}
\end{figure}

So the condition is easy to state and easy to check, with the caveat above kept
in view: \textbf{a chygraph is locally treelike only if its complexes meet in at
most one atom}, and on the random ensembles of Part~II that necessary condition
is sufficient as well. Uniform
hypergraphs drawn from a configuration model satisfy it with probability
approaching one. Multiplex and interacting networks satisfy it by construction.
Graphs decorated with triangles in the manner of \citet{newman2009} satisfy it
because the triangles are placed independently. What does not satisfy it is the
case one most wants: the maximal cliques of a real, geometrically clustered
graph, which overlap along edges as a matter of course
\citep{friedrich2015}.

This is stated in the chapter that defines the object rather than in a remark
at the end of the book. The next eleven chapters develop what the construction
buys; Chapter~\ref{ch:overlap} returns to Fig.~\ref{fig:price}b, measures the damage
on hyperbolic random graphs --- between $77$ and $100$ per cent of the answer
recovered --- and describes what would be needed to recover the rest. Clustering
is not one obstruction but a hierarchy of them: promoting
triangles removes the loops within triangles and exposes the loops between them.
Whether that hierarchy terminates usefully is, as of this writing, open.
